    \section{Motivation\label{ch1:motivation}}
    
\glsreset{AO}
Several reports, both classified and unclassified \cite{Jones2009,
ChiefScientistAirForce2010, Stone2022}, have identified \gls{AO} effects at
hypersonic speeds as a matter of critical importance to national security.
Applied research in this area is actively conducted not only at national
laboratories, such as \gls{LLNL}, \gls{JHU}-\gls{APL}, and \gls{LANL}, but
also by major \gls{DoD} contractors, including Lockheed Martin Corp.,
Northrop Grumman Corp., and Raytheon Technologies. As expected, much of this
applied work remains unavailable in the open literature. Despite these
efforts, a fundamental gap in understanding persists regarding the influence of
high-enthalpy flows at hypersonic speeds on \gls{AO} applications. \\

This dissertation seeks to address several existing knowledge gaps through
high-fidelity \gls{CFD} simulations, with the objective of characterizing the
influence of high-enthalpy flows on \gls{AO}-relevant \gls{EM}
properties, including the index of refraction, dielectric constant, electric
susceptibility, and polarizability. These properties are critical not only for
\gls{AO} applications but also for radar systems. While light beams and radar
signals are both \gls{EM} waves, they operate at vastly different frequency
ranges, leading to potentially distinct interactions with a high-enthalpy,
nonequilibrium medium. Understanding these interactions is essential for
predicting and mitigating performance degradation in optical and
radio-frequency sensing systems operating in extreme aerothermal
environments. \\

This work focus primarily on \gls{Lidar} sensor applications but it is also
relevant to \glspl{DEW}. The range of frequencies considered in this work is
constrained by available experimental data, and these limitations will be
specified where applicable.\\

        \Cref{fig:ch1strategicWeapon} illustrates the notional flight paths for
        various long-range weapons. Accurate target detection is of paramount
        importance in these scenarios. Cruise missiles generally travel at
        subsonic speeds, whereas ballistic missiles and boost-glide vehicles
        can reach speeds in the hypersonic regime. At these hypersonic
        velocities, the vehicles are subject to intense aerodynamic heating and
        dynamic pressure, resulting in complex hypersonic flow fields around
        the vehicle. These conditions can significantly degrade seeker
        performance, impacting sensors' ability to lock onto targets due to
        thermal interference, signal distortion and change in \gls{EM} properties. 
        Advanced guidance systems must compensate for these factors to maintain
        precision in target acquisition and engagement.
    \begin{figure}[H]
        \makebox[\textwidth][c]{
            \includegraphics[width=0.8\textwidth]{chapter1/strategicWeapons}}
        \caption{Notional flight paths of
        hypersonic boost-glide missiles, ballistic missiles, and cruise
        missiles. (\gls{CSBA} graphic)}\label{fig:ch1strategicWeapon}
    \end{figure}

        In addition, hypersonic vehicles can be detected by enemy ground-based
        radar systems. \Cref{fig:radarStrategicWeapons} illustrates the flight
        trajectories of an \gls{ICBM} and a \gls{HGV}. It is evident that
        \glspl{ICBM} are more easily detected compared to \glspl{HGV} primarily
        due to their higher apogee (the point of maximum altitude) during
        flight. \glspl{ICBM} typically travel through space along a ballistic
        arc, reaching altitudes that expose them to long-range radar systems
        for an extended period.
        %https://news.usni.org/2023/01/27/report-to-congress-on-hypersonic-missile-defense-2
        \pic{0.9}{chapter1/radarStrategicWeapons}{Hypersonic weapons}{fig:radarStrategicWeapons}

        In contrast, \glspl{HGV} operate at lower altitudes within the
        atmosphere, following a maneuverable, non-ballistic trajectory. This
        allows them to remain below the radar horizon for a longer duration,
        thereby reducing their detectability and increasing their
        survivability. However, because \glspl{HGV} fly within the denser
        layers of the atmosphere at extremely high velocities, they are subject
        to significant aerodynamic heating, turbulence, and real gas
        effects—all of which can degrade the performance of onboard optical
        sensors and seekers. \\

        Maintaining precise and accurate \gls{GNC} under such extreme
        conditions is critical to mission success. The ability of these systems
        to rapidly adapt to a dynamically changing environment, maintain sensor
        fidelity, and accurately strike their target despite \gls{AO} 
        distortions is essential for the operational effectiveness of
        hypersonic weapons.

        A crude approximation for the index of refraction in the atmosphere was proposed by Smith and Weintraub \cite{Smith1953}
        and it is given by 
        \begin{equation}\label{eq:crudeIndexApproximation}
            n(h) \approx 1 + \frac{K_1}{T(h)} \Parenthesis{ p(h) +
            K_2\frac{e(h)}{T(h)}}\,,
        \end{equation}
        where $K_1 =
        79\UnitsFrac{K}{mB}$, $K_2=4800\Units{K}$, and $e$ is the partial pressure
        of water vapor in millibars. \\

        \Cref{fig:ch1AtmosphericOptics} illustrates the index of refraction
        as a function of altitude for a \textit{blue-sky} scenario. The most
        significant change in the index of refraction occurs at approximately
        $11\Units{km}$ above sea level, within the tropopause, near the
        boundary between the troposphere and the stratosphere. This is also the
        typical cruising altitude for commercial aircraft and many high-speed
        jets. 
        \pic{0.8}{chapter1/atmosphericOptics}{Atmospheric index of
        refraction}{fig:ch1AtmosphericOptics}

        This region is critical for sensor accuracy because it is relatively
        close to the ground (where high-resolution sensing is often required)
        and is frequently associated with atmospheric turbulence. Both
        turbulence and variations in chemical composition can influence
        \gls{EM} wave propagation. In hypersonic flows, the chemical
        composition is not constant due to chemical nonequilibrium, resulting
        in additional spatial and temporal variations in the index of
        refraction. These variations can degrade the performance of optical
        systems, radar, and other \gls{EM}-based sensing technologies operating
        in or through this altitude band. \\

Active seekers are not engaged throughout the entire flight trajectory;
they are typically activated only during the terminal phase of flight.
This phase is generally characterized by lower speeds compared to midcourse
flight, although the exact conditions depend on the specific vehicle and
mission profile. The present work focuses on this region—the terminal
phase—because it is the portion of the trajectory where the vehicle is most
vulnerable to enemy defenses and interception. The specific conditions of the
terminal phase depend strongly on the flight trajectory and target objectives,
but are generally characterized by high speeds, intense aerodynamic heating
from atmospheric friction, and potential maneuvering to evade defenses. \\

        The index of refraction $\mathcal{N}$ consists of a real and complex
        component
        \begin{equation}\label{eq:totalIndexOfRefraction}
            \mathcal{N} = n + ik.
        \end{equation}
        The real
        part, $n$, governs the phenomenon of bending or refraction of
        light as it passes through a medium, while the complex part, $k$, is
        associated with attenuation or absorption of the wave as it
        propagates. The attenuation arises due to energy dissipation within the
        medium, which is particularly relevant in plasmas, where ionization
        effects contribute to optical losses. \\

        This work focus on analyzing the real part of
        \Cref{eq:totalIndexOfRefraction}, as it plays a crucial role in
        determining optical path differences and wavefront distortions in
        \gls{AO} applications. However, for a more complete picture of the role
        of the index of refraction,
        \Cref{fig:attenuationDB} presents the atmospheric attenuation at two
        different altitudes across a range of frequencies. The calculation of
        this attenuation requires a thorough understanding of the complex part
        of \Cref{eq:totalIndexOfRefraction}, as it governs the absorption and
        energy dissipation mechanisms in the medium.
        \pic{1.0}{chapter1/attenuationDB}{Attenuation at different frequencies
        \cite{ITU2022, Richards2010}}{fig:attenuationDB}

        Understanding the effects of the environment through which an optical
        beam must travel are imperative to accurate target detection.
        Unfortunately, most of the engineers working in \gls{Lidar} and Radars
        are mainly concerned with atmospheric environment effects such as rain, clouds,
        fog, and snow. However, they typically do not consider the impact of
        aerodynamic effects, such as shock-waves and boundary-layer
        interactions, on optical and radar beams. \\

        For subsonic
        flows, it might be a valid assumption to ignore the aerodynamic effects
        on the optical beam; however, for supersonic, and especially for
        hypersonic flows, this assumption should be revisited. During a hypersonic
        flight the temperatures behind the shock-waves are very high, which
        will not only induce thermal noise on the sensor but will also affect the
        electromagnetic properties that will degrade the sensor's
        performance leading to optical distortions and power loss. Being able to accurately compensate for
        these unpreventable effects is of extreme importance when designing
        sensors to detect the position and range of objects moving at
        high speeds. \\

        \Cref{fig:ch1AeroOpticsSchematic}, shows the path that a beam in a
        hypersonic flight will follow. The beam has to travel across the
        isolated gas surrounding the laser, the radome, the boundary layer, shock layer, and lastly
        atmospheric environment. Traditionally, engineers ignore aerodynamic
        effects such as shock-waves, and expansion-waves for their analysis.
        However, for the case
        of hypersonic flights, the aerodynamic effects cannot be ignored
        anymore because the electromagnetic properties across the boundary
        layer and shock layer drastically change. 
        \pic{0.8}{chapter1/aeroOpticsSchematic}{Optics Schematic}{fig:ch1AeroOpticsSchematic}

        Quadrants I and III correspond to low-flow resolution cases, such as
        laminar solvers, and \gls{RANS} simulations in which the turbulent fluctuations
        are not accounted for. \gls{AO} degradations, such as
        \gls{BSE}, \gls{OPL}, \gls{OPD}, and \gls{WD} require turbulent
        fluctuations to be accounted for. These quantities are defined
        in \cref{ch2:opticsOverview}.\\

        Quadrants II and IV corresponds to high-enthalpy resolution in which
        thermal and chemical nonequilibrium are accounted for. These
        simulations will
        provide a higher fidelity of flow properties, such as \gls{GD},
        polarizability, electric susceptibility, medium's permittivity, and
        index of refraction. These properties are discussed in more details in
        \cref{ch2:opticsOverview}.
        \begin{table}[ht!]
    \hspace*{-1cm}
    \begin{tikzpicture}
        % Draw the table
        \node (table) at (0,0) {
            \scalebox{0.75}{
                \begin{tabular}{|l|l|}
                    \hline
                    \textbf{Quadrant III:} High-Enthalpy, Low-Flow Resolution & \textbf{Quadrant IV:} High-Enthalpy, High-Flow Resolution \\
                    \makecell[l]{\textit{Physics captured:} chemical and thermal\\ nonequilibrium, shock structures, bulk unsteadiness,\\ bulk density gradients.} &
                    \makecell[l]{\textit{Physics captured:} chemical and thermal\\ nonequilibrium, shock structures, bulk unsteadiness,\\ bulk density gradients.} \\
                    \makecell[l]{\textit{AO quantities measured:} Gladstone-Dale Constant,\\ Polarizability, Dielectric Constant,\\ BSE, OPL, WD.} &  
                    \makecell[l]{\textit{AO quantities measured:} Gladstone-Dale Constant,\\ Polarizability, Dielectric Constant,\\ BSE, OPL, WD, OPD.} \\
                    \hline
                    \textbf{Quadrant I:} Low-Enthalpy, Low-Flow Resolution & \textbf{Quadrant II:} Low-Enthalpy, High-Flow Resolution \\ 
                    \makecell[l]{\textit{Physics captured:} shock structures, bulk\\ unsteadiness, bulk density gradients.} &
                    \makecell[l]{\textit{Physics captured:} shock structures, bulk\\ unsteadiness, bulk density gradients.} \\ 
                    \makecell[l]{\textit{AO quantities measured:} Gladstone-Dale Constant,\\ Polarizability, Dielectric Constant,\\ BSE, OPL, WD.} &
                    \makecell[l]{\textit{AO quantities measured:} Gladstone-Dale Constant,\\ Polarizability, Dielectric Constant,\\ BSE, OPL, WD, OPD.} \\
                    \hline
                \end{tabular}
            }
        };

        % Axes: Adjust coordinates based on final table scale
        \draw[->, thick] (-8,-3) -- (-8, 3) node[midway, rotate=90, above] {Increasing Flow Enthalpy};
        \draw[->, thick] (-8,-3) -- (2,-3) node[midway, below] {Increasing Flow Resolution};
    \end{tikzpicture}
    \caption{Schematic of the fidelity of aero-optical simulations.}
    \label{tab:ch1AeroOpticalFidelities}
\end{table}


\Cref{fig:gdFidelities} provides a summary of the different approaches
used to determine the \gls{GD} constants, which is used to calculate the index of
refraction. To the best of the author's knowledge,
most published work assumes a constant \gls{GD} across the flow. This may be a
reasonable assumption for low-enthalpy flows, but it becomes invalid in
high-enthalpy regimes. At low enthalpy, there are no chemical reactions that
alter the flow composition or excite internal energy modes. However, as the
temperature increases, chemical reactions occur that modify the composition of
the gas, motivating the need for species-specific \gls{GD} constants. One could
further extend this model to account for both internal state-specific and
species-specific \gls{GD} values.
        \begin{figure}[H]
            \makebox[\textwidth][c]{\includegraphics[width=0.9\textwidth]{chapter1/gdFidelities}}
            \caption{\acrlong{GD} fidelities \label{fig:gdFidelities}}
        \end{figure}

This work investigates \gls{AO} properties using the \gls{GD}
constant, which is evaluated at three different fidelity levels. These
fidelity levels are highlighted in orange in \Cref{fig:gdFidelities}. A
low-enthalpy, high-enthalpy, high-flow case (Quadrant II) is used to investigate
\textit{temperature-dependent \gls{GD}}, while a high-enthalpy, low-flow case
(Quadrant III) is used to investigate \textit{species-specific \gls{GD}}. A
\gls{StS} case is used to investigate \textit{internal-state-specific,
species-dependent \gls{GD}}.\\

        \Cref{fig:turbulenceFlows} illustrates how the flow field appears at increasing levels of simulation fidelity, from \gls{RANS} to \gls{LES} and finally to \gls{DNS} resolution for a jet flow. Each increase in \gls{CFD} resolution captures progressively finer turbulent structures, which in turn influence the predicted optical properties of the flow.
        \begin{figure}[H]
            \makebox[\textwidth][c]{\includegraphics[width=0.7\textwidth]{chapter1/turbulenceFlows}}
            \caption{Turbulence fidelities \cite{idealsimulations2025}\label{fig:turbulenceFlows}}
        \end{figure}



